\scanpage{25}
\diagram{\(\begin{array}{ccccc}
H^0&H^1&H^2&H^3&H^4\\
k&V&k\xi+W&\check V&k\xi^2
\end{array}\)}
\margin{Algèbres de Hodge d'ordre 2}

$V,W$
\[
\begin{aligned}
{\textstyle\bigwedge}^{2}V&\xrightarrow{\alpha}k\xi+W&
\alpha(x\wedge y)&=\varphi(x,y)\xi+\psi(x,y)&
\varphi(x,y)&\quad\text{forme alternée}\\
&&&&\psi(x,y)&\quad\text{$W$-\unsure{v|p}l. bil. alt. dans $W$}\\
\operatorname{Sym}^{2}(W)&\xrightarrow{\beta}k\xi^2&
\beta(xy)&=Q(x,y)\xi^2&
Q&\quad\text{forme bil. sym. sur $W$}
\end{aligned}
\]
\deleted{\unsure{Permet de recons}} La connaissance de $V,W,\varphi,\psi,\unsure{\bar Q|Q}$ permet de reconstituer $H^*$ à isom canonique près.
\[
\begin{aligned}
H^0\times H^i&\mathrel{\unsure{-|\to}}H^i:\quad\text{clair par}\quad 1.\unsure{x|\alpha}=x.1=x\qquad\text{tt }x.\\
H^1\times H^i&\longrightarrow H^{i+1}\quad\left|
\begin{aligned}
\deleted{\unsure{\text{d'où}|\text{dont}}}\ i=1\quad &xy=\varphi(x,y)\xi+\psi(x,y) &x,y\in V\\
i=2\quad &\left|\begin{aligned}
x\xi&=\xi x=\widetilde\varphi(x): &&\deleted{\unsure{(z)}}y\to\varphi(y,x)\\
xw&=wx=\rho(w,x):&&y\longrightarrow Q(\psi(y,x),w)
\end{aligned}\right. &\left|\begin{aligned}x&\in V\\y&\in V\\w&\in W\end{aligned}\right.\\
i=3\quad &xy'=-y'x=\deleted{\unsure{+}}\langle x,y'\rangle\xi^2 &x\in V,\ y'\in\check V\\
\deleted{i=4}\quad &\deleted{\unsure{x\xi^2=\xi}}
\end{aligned}\right.
\end{aligned}
\]
\[
H^2\times H^i\longrightarrow H^{i+2}\ \left|\quad j=2:\quad
\left|\begin{aligned}
\xi\cdot\xi&=\xi^2\\
\xi w=w\xi&=0\quad(w\in W)\\
ww'&=Q(w,w')\xi^2
\end{aligned}\right.\right.
\]
\margin{$(*)$ On pose $\widetilde\varphi:V\mathrel{\unsure{-|\to}}\check V$, avec $\langle x,\widetilde\varphi(y)\rangle=\varphi(x,y)$}

Quels que soient les choix de $\varphi,\psi$, on trouve une loi de composition anticommutative, admettant $1$ comme unité.

\uline{Condition d'associativité} s'écrit $(xy)(tz)=(tx)(yz)$, sous la forme
\[
\deleted{\unsure{\varphi}}\quad
\boxed{\varphi(x,y)\varphi(t,z)+Q(\psi(x,y),\psi(t,z))=\varphi(t,x)\varphi(y,z)+Q(\psi(t,x),\psi(y,z))}
\]
% slanted
\margin{exprime aussi que la forme $(xy|zt)$ en $xyzt$ est alternée en $\unsure{y,z|x,y,z|y^2}$ (donc en les \unsure{quatre lettres} !)}
\uline{Condition de Lefschetz}
\[
\begin{aligned}
\xi^{\unsure{2|?}}:H^0&\simeq H^4 &&\text{tjrs satisfait}\\
\xi:H^1&\longrightarrow H^3 &&\text{sss}\quad\boxed{\varphi\ \text{non dégénérée}}
\end{aligned}
\]
\uline{Condition de Poincaré}
\[
\begin{aligned}
H^4&\simeq k &&\text{tjrs satisfait}\\
H^0\times H^4&\longrightarrow H^4\ \text{et}\ H^1\deleted{\unsure{-}}\times H^3\mathrel{\unsure{-|\to}}H^4 &&\text{non dég : \uline{\unsure{itou|itrs|idem}}}\\
H^{\unsure{2|?}}\times H^2&\mathrel{-}H^{\unsure{4|2}} &&\text{non dég sss}\quad\boxed{Q\ \text{non dégénérée}}
\end{aligned}
\]
% archive sheet number: 24
